\documentclass[11pt]{article}
\usepackage{amsmath, amssymb, amsfonts, mathtools}
%\usepackage{hyperref}
\usepackage[margin=1in]{geometry}

\newcommand{\tr}{\mathrm{tr}}
\newcommand{\ii}{\mathrm{i}}
\newcommand{\dd}{\mathrm{d}}
\newcommand{\eps}{\varepsilon}
\newcommand{\mO}{\mathcal{O}}
\usepackage{
  jheppub
} % for details on the use of the package, please see the JINST-\
\newcommand{\nn}{\nonumber}
\usepackage{braket}
\newcommand{\lb}{\left(}
\newcommand{\rb}{\right)}
\usepackage{amsmath, amssymb}
\newcommand{\NL}[1]{\textcolor{blue}{#1}}

\begin{document}
  \section{A conjecture}

  Consider two energy eigenstates of a CFT $\ket{\mO_\alpha}$ and
  $\ket{\mO_\beta}$ corresponding to Virasoro primaries $\mO_{\alpha}$ and
  $\mO_{\beta}$ of dimensions $h_{\alpha}$ and $h_{\beta}$, respectively.
  Consider a subregion of angular size $A=(-\theta,\theta)$ on a spatial circle
  of circumference $2\pi$, with $\theta=\pi x$. We denote by $\rho_{\alpha}$ and
  $\rho_{\beta}$ the reduced states on $A$ in their corresponding vectors.

  The conjecture is that for the mixture $1/2(\rho_{\alpha}+\rho_{\beta})$ we
  have the following asymptotic behavior of relative entropy
  \begin{eqnarray}
    &&S(x)=S(\rho_\alpha|(\rho_\alpha+\rho_\beta)/2)\nn\\
    &&S(x)=A_0 x^{\delta_0}, \qquad S(1-x)=A_1(1-x)^{\delta_1},\nn\\
    &&\delta_1\leq \delta_0\ .
  \end{eqnarray}

  In particular, we will be interested in the special case of massless 2d free
  boson CFT and energy eigenstates created by the action of vertex operator
  states $V_{\alpha}=e^{i\alpha\phi}$ and $V_{\beta}= e^{i\beta\phi}$ at the
  center of radial quantization.

  \section{Relative entropy for mixed excited states}

  \subsection*{Setup and Replica Trick}

  Consider a conformal field theory (CFT) in $1+1$ dimensions on a cylinder
  $\mathcal{C}_{1}$ of circumference $R$. We focus on a spatial subsystem $A$ defined
  by the interval $(-l/2, l/2)$ at a fixed time slice. We simplify our notation
  by taking $A=(-\theta,\theta)$ on a cylinder of circumference $2\pi$ such that
  $\theta=\pi x$.

  Our states of interest are the reduced density matrices corresponding to
  global excited coherent states. These states are generated by the action of vertex
  operators $V_{\alpha}=:e^{i\alpha \phi}:$ on the vacuum at past infinity in imaginary
  time on the Euclidean cylinder $\mathcal{C}_{1}$:
  \begin{equation}
    |\alpha\rangle = \lim_{u\to i\infty}\frac{V_{\alpha}(u)|\Omega\rangle}{\braket{V_{-\alpha}(-u)V_\alpha(u)}^{1/2}}
  \end{equation}
  The reduced density matrix of the state $|\alpha\rangle$ on the subsystem $A$
  is denoted by $\rho_{\alpha}$. A simple way to keep track of normalization of the
  density matrix to defined normalized operators
  \begin{eqnarray}
    \tilde{V}_{\alpha}(u^+)\tilde{V}_{-\alpha}(u^-)=\frac{V_{\alpha}(u^{+})V_{-\alpha}(u^{-})}{\braket{V_{\alpha}(u^+)V_{-\alpha}(u^-)}_{\mathcal{C}_1}}
  \end{eqnarray}

  We aim to compute the relative entropy $S(\rho || \sigma_{y})$ where
  $\rho = \rho_{\alpha}$ and the reference state is the mixture
  $\sigma_{y}= y\rho_{\alpha}+(1-y) \rho_{\beta}$ for $0\leq y\leq 1$. The standard
  replica trick for relative entropy relates this quantity to the analytic continuation
  of trace functionals that break replica symmetry. Specifically, for the mixed state
  case, we utilize the following identity derived from the generalized replica
  trick for normalized density matrices:
  \begin{eqnarray}
    S\left(\rho_\alpha \Big\| y\rho_\alpha+(1-y)\rho_\beta\right)&&= \partial_{n}
    \log \left[ \frac{\text{tr}(\rho_{\alpha}^{n})}{\text{tr}(\rho_{\alpha}(y\rho_{\alpha}+(1-y)\rho_{\beta})^{n-1})}
    \right]_{n\to 1}\nn\\
    && =\partial_n\log\left[
    \frac{\braket{\prod_{i=1}^{n} \tilde{V}_{\alpha}(u^+_i)\tilde{V}_{-\alpha}(u^-_i)}_{\mathcal{C}_n}}{\braket{\tilde{V}_\alpha(u^+_n)\tilde{V}_{-\alpha}(u^-_n)\mO^{(n-1)}_{\alpha|\beta}}_{\mathcal{C}_n}}\right]_{n=1}
  \end{eqnarray}
  where
  \begin{eqnarray}
    &&\mO_{\alpha|\beta}^{(n-1)}=\sum_{p=0}^{n-1}y^{n-1-p}(1-y)^p\sum_{\substack{w\in\{\alpha,\beta\}^{n-1}\\ N_\beta(w)=p}}
    \tilde{V}_{w_1}(u_1^+)\tilde{V}_{-w_1}(u_1^-)\cdots \tilde{V}_{w_{n-1}}(u_{n-1}^+)\tilde{V}_{-w_{n-1}}(u_{n-1}^-)\nn\\
  \end{eqnarray}
  and $w=(w_{1},\dots,w_{n-1})$ is an ordered word and $N_{\beta}(w)$ counts how
  many letters equal $\beta$, and the points are on the $n$-cover of cylinder
  $\mathcal{C}_{n}$ at $\pm i\infty$ on each sheet.

  The conformal map
  \begin{eqnarray}
    z=\lb\frac{\sin((u-\theta)/2)}{\sin((u+\theta)/2)}\rb^{1/n}
  \end{eqnarray}
  sends the interval $A=(-\theta,\theta)$ to $A=(0,\infty)$ unifomizing
  $\mathcal{C}_{n}$ to the complex plane. Under this transformation Virasoro primaries
  transform as
  \begin{eqnarray}
    &&\mO(u,\bar{u})=\lb\frac{\partial z}{\partial u} \rb^{h_\alpha}\lb\frac{\partial
    \bar{z}}{\partial\bar{u}}
    \rb^{\bar{h}_\alpha}\mO(z,\bar{z})\nn\\
    &&V_\alpha(u_k,\bar{u}_k)=\lb\frac{z_{k,n}\Lambda}{n}\rb^{h_\alpha}\lb\frac{\bar{z}_{k,n}\Lambda}{n}\rb^{\bar{h}_\alpha}V_\alpha(z_{k,n},\bar{z}_{k,n}),
  \end{eqnarray}
  where in our example
  \begin{eqnarray}
    z^\pm_{k,n}=e^{\frac{i\pi}{n}(2k\pm x)},\qquad \Lambda=\frac{\sin(\theta)}{\cos(\theta)-\cos(u)}
  \end{eqnarray}
  Therefore, for normalized operators $\tilde{V}$ the contribution of $\Lambda$
  terms cancel and we find
  \begin{eqnarray}
    \tilde{V}_\alpha(u_k,\bar{u}_k)&&=\lb\frac{z_{k,n}}{n z_{1,1}}\rb^{h_\alpha}\lb\frac{\bar{z}_{k,n}}{n
    z_{1,1}}\rb^{\bar{h}_\alpha}\tilde{V}_\alpha(z_{k,n},\bar{z}_{k,n})\nn\\
    &&=\lb\frac{e^{i\pi(2k\pm x)(1-n)/n}}{n }\rb^{h_\alpha}\lb\frac{e^{-i\pi(2k\pm
    x)(1-n)/n}}{n}\rb^{\bar{h}_\alpha}\tilde{V}_\alpha(z_{k,n},\bar{z}_{k,n})
  \end{eqnarray}
  For scalar (spinless) operators in $2d$ we have $h_{\alpha}=\bar{h}_{\alpha}$
  and we find
  \begin{equation}
    \tilde{V}_{\alpha}(u_{k},\bar{u}_{k})&&=n^{-2h_\alpha}\tilde{V}_{\alpha}(z_{k,n}
    ,\bar{z}_{k,n})
  \end{equation}
  Further assuming that $h_{\alpha}=h_{-\alpha}$ the relative entropy replica
  trick in terms of the correlator on the complex plane gives
  \begin{eqnarray}
    &&S(\rho_\alpha\|y\rho_\alpha+(1-y)\rho_\beta)=\partial_n\log\mathcal{F}_n\big|_{n\to 1}\nn\\
    &&\mathcal{F}_n\equiv \left[
    \frac{\braket{\prod_{k=1}^{n} \tilde{V}_{\alpha}(z^+_{k,n})\tilde{V}_{-\alpha}(z^-_{k,n})}}{\braket{\tilde{V}_\alpha(z^+_{n,n})\tilde{V}_{-\alpha}(z^-_{n,n})\tilde{\mO}^{(n-1)}_{\alpha|\beta}}}\right]_{n=1}\nn\\
    &&\tilde{\mO}_{\alpha|\beta}^{(n-1)}=\sum_{p=0}^{n-1}y^{n-1-p}(1-y)^p\: n^{4p(h_\alpha-h_\beta)}
    \sum_{\substack{w\in\{\alpha,\beta\}^{n-1}\\ N_\beta(w)=p}} \prod_{k=1}^{n-1}\frac{V_{w_k}(z_{k,n}^{+})V_{-w_k}(z_{k,n}^{-})}{\braket{V_{\omega_k}(z^+_{1,1})V_{-\omega_k}(z^-_{1,1})}}\nn\\
  \end{eqnarray}
  % For operators $h_\alpha=\bar{h}_\alpha$, and $h_\beta=\bar{h}_\beta$, we have
  % \begin{eqnarray}
  %   \tilde{\mO}_{\alpha|\beta}^{(n-1)}=\sum_{p=0}^{n-1}y^{n-1-p}(1-y)^p\lb \frac{\Lambda^2}{n^2}\rb^{2p(h_\beta-h_\alpha)}\sum_{\substack{w\in\{\alpha,\beta\}^{n-1}\\ N_\beta(w)=p}}
  %   \prod_{k=1}^{n-1} \frac{V_{w_k}(z_k^+)V_{-w_k}(z_k^-)}{\braket{V_{w_k}(z^+)V_{-w_k}(z^-)}}
  % \end{eqnarray}

  \subsection*{Small $x$ limit}

  To simplify our notation we define $C_{\omega}=C^{L}_{\omega(-\omega)}$ for
  $\omega=\alpha,\beta$. In the small $x$ limit, we can expand in the OPE
  channel
  \begin{eqnarray}
    \frac{V_{w_k}(z_{k}^{+})V_{-w_k}(z_{k}^{-})}{\braket{V_{w_k}(z^+)V_{-w_k}(z^-)}}=\lb
    \frac{\sin(\pi x)}{\sin(\pi x/n)}\rb^{2h_{\omega_k}}\lb 1+ C_\omega(2\sin(\pi
    x/n))^{\Delta_L}\mO_L(e^{2\pi i k/n})+\cdots\rb\nn
  \end{eqnarray}
  where $\mO_{L}$ is the lightest Virasoro primary that we assume to have
  dimensions $(h_{L},\bar{h}_{L})$, weight $\Delta_{L}=h_{L}+\bar{h}_{L}$ less than
  two and spin $s_{L}=h_{L}-\bar{h}_{L}$ \footnote{Otherwise the lowest order
  term would come from stress tensor.}. Therefore, we find
  \begin{align}
    \label{eq:F_n}\mathcal{F}_{n} & =\frac{\,1+(2\sin(\pi x/n))^{2\Delta_L}f(n)+\cdots\,}{\displaystyle \sum_{p=0}^{\,n-1}y^{\,n-1-p}(1-y)^{p}\left(\frac{n^2\sin(\pi x/n)}{\sin(\pi x)}\right)^{4p(h_\alpha-h_\beta)}\sum_{\substack{w\in\{\alpha,\beta\}^{\,n-1}\\ N_\beta(w)=p}}\left[\,1+(2\sin(\pi x/n))^{2\Delta_L}g_{w}(n)+\cdots\,\right] }
  \end{align}
  where the leading correction $f(n)$ and $g_{p}(n)$ are given by the following sums
  over OPE coefficients and light operator correlators on the replica manifold:
  \begin{align}
    f(n) & =C_{\alpha}^{2}\sum_{k\neq k'=1}^{n}\langle \mO_{L}(e^{2\pi i k/n})\mO_{L}(e^{2\pi i k'/n})\rangle \nn    \\
         & =C_{\alpha}^{2}\frac{n}{2}\sum_{k=1}^{n-1}\frac{e^{-2\pi i s_L(k+k')/n}}{(2\sin(\pi k/n))^{2\Delta_L}}\ .
  \end{align}
  Rewrite the sum above by fixing the separation $m:=k-k' \ (\mathrm{mod}\ n)$. For
  each $m\in\{1,2,\dots,n-1\}$, set $k'=k-m$ (mod $n$). Then
  \[
    e^{-2\pi i s_L (k+k')/n}= e^{-2\pi i s_L (2k-m)/n}= e^{+2\pi i s_L m/n}\,e^{-4\pi
    i s_L k/n},
  \]
  and the denominator depends only on $m$. Hence the sum factorizes:
  \begin{equation}
    \sum_{m=1}^{n-1}\frac{e^{2\pi i s_L m/n}}{\bigl(2\sin\frac{\pi m}{n}\bigr)^{2\Delta_L}}
    \underbrace{\sum_{k=0}^{n-1} e^{-4\pi i s_L k/n}}_{=:G(n)}.
  \end{equation}
  The inner sum $G(n)$ is a geometric series with ratio $q=e^{-4\pi i s_L/n}$:
  \[
    G(n)=\sum_{k=0}^{n-1}q^{k}=
    \begin{cases}
      n, & q=1,     \\
      0, & q\neq 1.
    \end{cases}
  \]
  Therefore $S(n)=0$ for integer $n$ unless $q=1$, i.e.
  \[
    e^{-4\pi i s_L/n}=1 \quad\Longleftrightarrow\quad \frac{4s_{L}}{n}\in\mathbb{Z}
    .
  \]
  Now, since $4s_{L}/n$ is not an integer, the numerator $e^{2\pi i s_L m/n}$ is
  at most a sign. For $s_{k}=0$, we are left with
  \begin{align}
     & S_{2}(n)=\frac{n}{2}\sum_{k=1}^{n-1}(2\sin(\pi k/n))^{-2\Delta_L},\qquad S_{2}(1)=0,
  \end{align}
  whose analytic continuation of $S_{2}(n)$ was discussed in
  \cite{calabrese2010entanglement}
  \begin{align}
    S'_{2}(1)=\frac{\sqrt{\pi}\Gamma(\Delta_{L}+1)}{4\Gamma(\Delta_{L}+3/2)}\ .
  \end{align}
  In a general CFT, there could be several inequivalent choices of the lightest primary
  operator $\mO_{L}$ that contribute at the same order. In this case, the leading
  correction $f(n)$ would be given by the sum over all such operators, that
  following \cite{calabrese2010entanglement} we denote by
  \begin{align}
    f(n) & =\mathcal{N}S_{2}(n)
  \end{align}
  with $\mathcal{N}=2$ for free boson and $\mathcal{N}=1$ for Ising model. For similicity,
  we set $\mathcal{N}=1$ in the remainder of this calculation.

  Similarly, we have
  \begin{align}
    g_{w}(n)= & C_{\alpha}\sum_{k=1}^{n-1}\frac{C_{\omega_k}\mathcal{N}}{(2\sin(\pi k/n))^{2\Delta_L}}+\sum_{k\neq k'=1}^{n-1}\frac{C_{\omega_k}C_{\omega_{k'}}}{(2\sin(\pi |k-k'|/n))^{2\Delta_L}}, % \nn\\
    % &=C_\alpha C_\omega \frac{2 S_2(n)}{n}+\sum_{k\neq k'=1}^{n-1}
    % \frac{C_{\omega_k}C_{\omega_{k'}}}{(2\sin(\pi |k-k'|/n))^{4h_L}}
  \end{align}

  We can simplify the denominator of $\mathcal{F}_{n}$ by noting that
  \begin{align}
    % &\sum_{\substack{w\in\{\alpha,\beta\}^{\,n-1}\\
    %                        N_\beta(w)=p}} 1={n-1 \choose p},\nn\\
     & \sum_{\substack{w\in\{\alpha,\beta\}^{\,n-1}\\ N_\beta(w)=p}}\left[\,1+(2\sin(\pi x/n))^{2\Delta_L}g_{w}(n)+\cdots\,\right]={n-1 \choose p}+(2\sin(\pi x/n))^{2\Delta_L}\bar{g}_{p}(n)+\cdots\,\nn                                   \\
     & \bar{g}_{p}(n) =\sum_{\substack{w\in\{\alpha,\beta\}^{\,n-1}\\ N_\beta(w)=p}}g_{w}(n)=C_{\alpha}\sum_{k=1}^{n-1}\frac{A_{k}}{(2\sin(\pi k/n))^{2\Delta_L}}+\sum_{k\neq k'=1}^{n-1}\frac{B_{kk'}}{(2\sin(\pi |k-k'|/n))^{2\Delta_L}},
  \end{align}
  where we have exchanged the order of the sum over words and the sum over
  replica indices and defined
  \begin{align}
     & A_{k}(p,n)=\sum_{\substack{w\in\{\alpha,\beta\}^{\,n-1}\\ N_\beta(w)=p}}C_{\omega_k}={n-2 \choose p}C_{\alpha}+{n-2 \choose p-1}C_{\beta}\equiv A(p,n)\nn                                                              \\
     & B_{kk'}(p,n)=\sum_{\substack{w\in\{\alpha,\beta\}^{\,n-1}\\ N_\beta(w)=p}}C_{\omega_k}C_{\omega_{k'}}={n-3 \choose p}C_{\alpha}^{2}+2{n-3 \choose p-1}C_{\alpha}C_{\beta}+{n-3 \choose p-2}C_{\beta}^{2}\equiv B(p,n), % \nn\\
    % &=\frac{(n-1-p)(n-2-p)}{(n-1)(n-2)}C_{\alpha}^2+\frac{2p(n-1-p)}{(n-1)(n-2)}C_{\alpha}C_{\beta}+\frac{p(p-1)}{(n-1)(n-2)}C_{\beta}^2\nn
  \end{align}
  where we performed the sums over words for fixed $k\neq k'$ and found that the
  result is independent of $k$ and $k'$. Therefore, we obtain
  \begin{align}
    g_{p}(n) & =C_{\alpha}A(p,n)\sum_{k=1}^{n-1}\frac{1}{(2\sin(\pi k/n))^{2\Delta_L}}+B(p,n)\sum_{k\neq k'=1}^{n-1}\frac{1}{(2\sin(\pi |k-k'|/n))^{2\Delta_L}}\nn      \\
             & =(C_{\alpha}A(p,n)-B(p,n))\sum_{k=1}^{n}\frac{1}{(2\sin(\pi k/n))^{2\Delta_L}}+B(p,n)\sum_{k\neq k'=1}^{n}\frac{1}{(2\sin(\pi |k-k'|/n))^{2\Delta_L}}\nn \\
             & =\frac{2S_{2}(n)}{n}(C_{\alpha}A(p,n)-B(p,n))+B(p,n)S_{2}(n)\nn                                                                                          % \\
    %& =A(p,n)\frac{2S_{2}(n)}{n}+B(p,n)S_{2}(n)
  \end{align}
  and as a result we can write
  \begin{align}
     & \left[\,1+(2\sin(\pi x/n))^{2\Delta_L}g_{w}(n)+\cdots\,\right]={n-1 \choose p}+(2\sin(\pi x/n))^{2\Delta_L}g_{p}(n)+\cdots\,
  \end{align}

  Defining
  \begin{equation}
    D_{n}=\left(\frac{n^{2}\sin(\pi x/n)}{\sin(\pi x)}\right)^{2(h_\alpha-h_\beta)}
    ,\qquad A_{1}=1,
  \end{equation}
  we note the following identities
  \begin{align}
     & Z(n)\equiv \sum_{p=0}^{\,n-1}y^{\,n-1-p}((1-y)D_{n})^{p}{n-1\choose p}= (y+(1-y) D_{n})^{n-1}                                 \\
     & A_{\alpha}(n)\equiv \sum_{p=0}^{\,n-1}y^{\,n-1-p}((1-y)D_{n})^{p}{n-2\choose p}=y(y+(1-y) D_{n})^{n-2}                        \\
     & A_{\beta}(n)\equiv \sum_{p=0}^{\,n-1}y^{\,n-1-p}((1-y)D_{n})^{p}{n-2\choose p-1}=(1-y)D_{n}(y+(1-y) D_{n})^{n-2}              \\
     & B_{\alpha\alpha}(n)\equiv\sum_{p=0}^{\,n-1}y^{\,n-1-p}((1-y)D_{n})^{p}{n-3\choose p}= y^{2}(y+(1-y)D_{n})^{n-3}               \\
     & B_{\alpha\beta}(n)\equiv\sum_{p=0}^{\,n-1}y^{\,n-1-p}((1-y)D_{n})^{p}{n-3\choose p-1}= y(1-y) D_{n}(y+(1-y)D_{n})^{n-3}       \\
     & B_{\beta\beta}(n)\equiv\sum_{p=0}^{\,n-1}y^{\,n-1-p}((1-y)D_{n})^{p}{n-3\choose p-2}= (1-y)^{2}D_{n}^{2}(y+(1-y)D_{n})^{n-3},
  \end{align}
  For simplicity, we define
  \begin{align}
     & X(n)=(y+(1-y)D_{n}), \qquad X(1)=1\nn                 \\
     & Y_{\alpha\beta}(n)=(y C_{\alpha}+(1-y)D_{n}C_{\beta})
  \end{align}
  so that
  \begin{align}
    B(p,n) & =X(n)^{n-3}\lb C_{\alpha}^{2}B_{\alpha\alpha}(n)+2C_{\alpha}C_{\beta}B_{\alpha\beta}(n)+C_{\beta}^{2}B_{\beta\beta}(n)\rb\nn \\
           & =Y_{\alpha\beta}(n)^{2}X(n)^{n-3}\nn                                                                                         \\
    A(p,n) & = C_{\alpha}A_{\alpha}(n)+C_{\beta}A_{\beta}(n)=Y_{\alpha\beta}(n)X(n)^{n-2}
  \end{align}
  Then, in the denominator we have
  \begin{align}
     & X(n)^{n-1}\lb 1+\lb\frac{2\pi x}{n}\rb^{2\Delta_L}\Delta(n)+\cdots \rb
  \end{align}
  where
  \begin{align}
    \Delta(n) & \equiv \frac{2S_{2}(n)}{n}\frac{(C_{\alpha}A(p,n)-B(p,n))}{X(n)^{n-1}}+\frac{B(p,n)}{X(n)^{n-1}}S_{2}(n)\nn                                                          \\
              & =\frac{2S_{2}(n)}{n}\lb \frac{C_{\alpha}Y_{\alpha\beta}(n)}{ X(n)}-\frac{Y_{\alpha\beta}(n)^{2}}{X(n)^{2}}\rb+S_{2}(n)\lb \frac{Y_{\alpha\beta}(n)^{2}}{X(n)^{2}}\rb \\
              & =S_{2}(n)\lb \frac{Y_{\alpha\beta}(n)}{X(n)}\left(\frac{2C_{\alpha}}{n}\right)+\frac{Y_{\alpha\beta}^{2}(n)}{X(n)^{2}}\lb 1-\frac{2}{n}\rb\rb                        %& g_{1}(n)=\frac{2(y C_{\alpha}+(1-y)D_{n}C_{\beta})}{n(y+(1-y) D_{n})}+\frac{(yC_{\alpha}+(1-y)D_{n}C_{\beta})^{2}}{(y+(1-y) D_{n})^{2}}
  \end{align}
  Putting this back into (\ref{eq:F_n}) we find
  \begin{align}
    \log\mathcal{F}_{n}=(1-n)\log X(n)+\log\left[\frac{1+\lb \frac{2\pi x}{n}\rb^{2\Delta_L}S_{2}(n)\lb C_{\alpha}^{2}+\cdots\rb}{1+\lb \frac{2\pi x}{n}\rb^{2\Delta_L}\Delta(n)+\cdots}\right]
  \end{align}
  We first note that taking the $n$-derivative of the first factor in $\mathcal{F}
  _{n}$ and sending $n\to 1$ vanishes
  \begin{align}
    \partial_{n}\left[(1-n)\log(y+(1-y)D_{n})\right]\Big|_{n=1}=0
  \end{align}
  and we are left with
  \begin{align}
    S(\rho_{\alpha}\|y\rho_{\alpha}+(1-y)\rho_{\beta}) & =\partial_{n}\left[\lb \frac{2\pi x}{n}\rb^{2\Delta_L}(S_{2}(n) C_{\alpha}^{2}-\Delta(n))+\cdots\right]_{n=1}
  \end{align}
  Since $S_{2}(1)=0$ and $X(1)=1$, we have
  \begin{align}
    S(\rho_{\alpha}\|y\rho_{\alpha}+(1-y)\rho_{\beta}) & = (2\pi x)^{2\Delta_L}S_{2}'(1)(C_{\alpha}^{2}-2C_{\alpha}Y_{\alpha\beta}(1)+Y_{\alpha\beta}(1)^{2})+\cdots\nn                \\
                                                       & =(2\pi x)^{2\Delta_L}\frac{\sqrt{\pi}\Gamma(\Delta_{L}+1)}{4\Gamma(\Delta_{L}+3/2)}(1-y)^{2}(C_{\alpha}-C_{\beta})^{2}+\cdots
  \end{align}
  where in the second line we have used $S_{2}(1)=0$. In the special case,
  $\alpha =\beta$ or $y\to 1$ the term vanishes as expected. In the case of vertex
  operators in free massless bosons, we have
  \begin{align}
    V_{\omega}(z,\bar{z})V_{-\omega}(0,0)=|z|^{-2\Delta_\omega}\lb 1+i\omega\lb z J(0)+\bar{z}\bar{J}(0)\rb+\omega^{2} (z\bar{z}J(0)\bar{J}(0)+\cdots)\rb
  \end{align}

  % \NL{HERE}
  % \begin{align}
  %    & \sum_{p=0}^{n-1}y^{n-1-p}((1-y)A)^{p}S(p)={}\nn                                                                                              \\
  %    & \frac{(n-1)(n-2)}{(n-1)(n-2)}C_{\alpha}^{2}y^{2}(y+(1-y)A)^{n-3}+\frac{2(n-1)(n-1)}{(n-1)(n-2)}C_{\alpha}C_{\beta}A y^{2}(y+(1-y)A)^{n-3}\nn \\
  %    & =(y+(1-y)A)^{n-3}(y C_{\alpha}+(1-y)A C_{\beta})^{2}
  % \end{align}
  % and the denominator of $\mathcal{F}_{n}$ becomes
  % \begin{align}
  %    & (y+(1-y)A)^{n-1}\left[1+(2\sin(\pi x/n))^{4h_L}h(n)+\cdots\right] \nn \\
  %    & h(n)=\lb\frac{y C_{\alpha}+(1-y)A_{n}C_{\beta}}{y+(1-y)A_{n}}\rb^{2}
  % \end{align}

  % Therefore, $\mathcal{F}_{n}$ in the small $x$ expansion becomes
  % \begin{align}
  %   \mathcal{F}_{n}= & \lb y+(1-y)A_{n}\rb^{1-n}\lb 1+(2\sin(\pi x/n))^{4h_L}S_{2}(n)(C_{\alpha}^{2}-h(n))+\cdots\rb\nn                                                      \\
  %   =                & \lb y+(1-y)A_{n}\rb^{1-n}\lb 1+F^{(1)}_{n}+\cdots\rb\nn                                                                                               \\
  %                    & F_{n}^{(1)}=(2\sin(\pi x/n))^{4h_L}S_{2}(n)\frac{A_{n}(C_{\alpha}-C_{\beta})(1-y)(A_{n}(C_{\alpha}+C_{\beta})(1-y)+2yC_{\alpha})}{(y+(1-y)A_{n})^{2}}
  % \end{align}

  % Taking the logarithm and $n$ derivative at $n\to 1$ of the first factor in $\mathcal{F}
  % _{n}$ vanishes
  % \begin{align}
  %   \partial_{n}\lb (1-n)\log(y+(1-y)A_{n})\rb_{n\to 1}=-\log(y+(1-y)A_{1})=0\ .
  % \end{align}
  % Therefore, the leading contribution to the relative entropy in the small $x$
  % limit comes from the second factor and since $S_{2}(1)=0$ we find
  % \begin{align}
  %   \partial_{n}\log\mathcal{F}_{n}\big|_{n\to 1}=\partial_{n}F_{n}^{(1)}\big|_{n\to 1}=(2\sin(\pi x))^{4h_L}S_{2}'(1)(1-y)(C_{\alpha}-C_{\beta})(C_{\alpha}+C_{\beta}+y(C_{\alpha}-C_{\beta}))\ .
  % \end{align}
  % Putting it all together, we find the leading small $x$ behavior of the relative
  % entropy
  % \begin{align}
  %   S(\rho_{\alpha}|y\rho_{\alpha}+(1-y)\rho_{\beta})=(2\sin(\pi x))^{4h_L}\frac{\sqrt{\pi}\Gamma(2h_{L}+1)}{4\Gamma(2h_{L}+3/2)}(1-y)(C_{\alpha}-C_{\beta})(C_{\alpha}+C_{\beta}+y(C_{\alpha}-C_{\beta}))+\cdots
  % \end{align}
  % For the special case $y=1/2$ and $\alpha=0$ we find
  % \begin{align}
  %   S(\rho_{0}|(\rho_{0}+\rho_{\beta})/2)=(2\sin(\pi x))^{4h_L}\frac{\sqrt{\pi}\Gamma(2h_{L}+1)}{8\Gamma(2h_{L}+3/2)}C_{\beta}^{2}+\cdots
  % \end{align}
  % In the special case of vertex operators in the massless free boson, the lowest
  % dimension primary is $\partial\phi$ and we find
  % \begin{align}
  %   S(\rho_{0}|(\rho_{0}+\rho_{\beta})/2)=(2\sin(\pi x))^{\beta^2}\frac{\sqrt{\pi}\Gamma(\beta^{2}+1)}{8\Gamma(\beta^{2}+3/2)}+\cdots
  % \end{align}

  \subsection*{Small $1-x$ limit}

  The final answer will be (\ref{eq:S-final}). To start simple, we take $\alpha=0$
  so that $\rho_{\alpha}$ is the vacuum state reduced to the subsystem.

  % In this case, we have
  % \begin{eqnarray}
  %   &&S(\rho_0\|y\rho_0+(1-y)\rho_\beta)=\partial_n\left[ \frac{(y\braket{1}_{\mathcal{C}_1}+(1-y)\braket{V_\beta(u^+) V_{-\beta}(u^-)}_{\mathcal{C}_1})^{n-1}}{\braket{\mO^{(n-1)}_\beta}_{\mathcal{C}_n}}\right]_{n=1}\nn\\
  %   &&\mO_\beta^{(n-1)}=\sum_{p=0}^{n-1}y^{n-1-p}(1-y)^p\sum_{\substack{w\in\{0,\beta\}^{n-1}\\ N_\beta(w)=p}}
  % V_{w_1}(u_1^+)V_{-w_1}(u_1^-)\cdots V_{w_{n-1}}(u_{n-1}^+)V_{-w_{n-1}}(u_{n-1}^-)\nn\\
  % \end{eqnarray}

  The most non-trivial component of this calculation is the term
  $\text{tr}\!\left (\rho_{\alpha}\,(y\rho_{\alpha}+(1-y)\rho_{\beta})^{n-1}\right
  )$. Since the operators do not commute in general, one should expand the \emph{power}
  as a sum over all words of length $n-1$ in the alphabet $\{\rho_{\alpha},\rho_{\beta}
  \}$ rather than using a naive binomial theorem. In the case $\alpha=0$ we
  write $\rho_{0}$ for the vacuum reduced density matrix and obtain the precise
  finite-$n$ identity
  \begin{equation}
    \label{eq:word-expansion-alpha0}\tr\!\left(\rho_{0}\,(y\rho_{0}+(1-y)\rho_{\beta}
    )^{n-1}\right) = \sum_{p=0}^{n-1}y^{\,n-1-p}(1-y)^{p}\sum_{\substack{w\in\{0,\beta\}^{n-1}\\ N_\beta(w)=p}}
    \tr\!\left(\rho_{0}\,\rho_{w_1}\rho_{w_2}\cdots \rho_{w_{n-1}}\right),
  \end{equation}
  where $w=(w_{1},\dots,w_{n-1})$ is an ordered word and $N_{\beta}(w)$ counts
  how many letters equal $\beta$.

  Equivalently, one may separate the ``naive binomial'' term (in which all
  $\rho_{\beta}$ are taken to sit to the right) plus an explicit remainder built
  from commutators. A convenient way to package this remainder is via the
  adjoint action $\mathrm{ad}_{\rho_0}(X)\equiv[\rho_{0},X]$:
  \begin{equation}
    \label{eq:binomial-plus-comm-alpha0}\tr\!\left(\rho_{0}\,(y\rho_{0}+(1-y)\rho
    _{\beta})^{n-1}\right) = \sum_{p=0}^{n-1}\binom{n-1}{p}y^{\,n-1-p}(1-y)^{p}\,
    \tr\!\left(\rho_{0}^{\,n-p}\rho_{\beta}^{\,p}\right) +\mathcal{C}_{n}(y),
  \end{equation}
  with the commutator correction
  \begin{equation}
    \label{eq:Cn-def}\mathcal{C}_{n}(y) \equiv \sum_{p=0}^{n-1}y^{\,n-1-p}(1-y)^{p}
    \sum_{\substack{w\in\{0,\beta\}^{n-1}\\ N_\beta(w)=p}}\tr\!\left(\rho_{0}\,\Delta
    _{w}\right),\qquad \Delta_{w}\equiv \rho_{w_1}\cdots\rho_{w_{n-1}}-\rho_{0}^{\,n-1-p}
    \rho_{\beta}^{\,p}.
  \end{equation}
  Each $\Delta_{w}$ can be reduced to a finite sum of nested commutators by repeatedly
  commuting $\rho_{0}$ past $\rho_{\beta}$ (so $\mathcal{C}_{n}(y)$ vanishes
  whenever $[\rho_{0},\rho_{\beta}]=0$; in particular, for the free boson vertex-operator
  states considered below the ordering dependence drops out).

  Using the operator-state correspondence, the trace
  $\text{tr}(\rho_{\alpha}^{n-p}\rho_{\beta}^{p})$ corresponds to a path
  integral on an $n$-sheeted Riemann surface with specific operator insertions. In
  the limit $n \to 1$, we can identify this with the correlation function of $p$
  operators of type $\beta$ and $q=n-p$ operators of type $\alpha$:
  \begin{equation}
    \label{fullcorr}C_{q,p}\equiv \langle \mathcal{O}_{\alpha}^{(q)}\mathcal{O}_{\beta}
    ^{(p)}\rangle
  \end{equation}
  As was pointed out, in the case of free bosons, the order of operators will not
  matter, and the commutator terms vanish.

  \subsection*{Operator Insertions on the Strip}

  We perform the calculation on a strip of width $2\pi$, obtained via the
  standard conformal mapping from the $n$-sheeted cylinder. The replica sheets are
  indexed by $k = 0, \dots, n-1$.

  The operator insertion points $t_{k}^{\pm}$ on the strip are:
  \begin{equation}
    t_{k}^{\pm}= \frac{2\pi k}{n}\pm \frac{\pi x}{n}
  \end{equation}
  where $x = l/R$ is the dimensionless interval size. For the correlator $\langle
  \mathcal{O}_{\alpha}^{(q)}\mathcal{O}_{\beta}^{(p)}\rangle$ with $q=n-p$, the
  operators are explicitly distributed as follows: On sheets $k=0, \dots, q-1$ ($\alpha$-sector)
  we insert $V_{\alpha}$ at $t_{k}^{+}$ and $V_{-\alpha}$ at $t_{k}^{-}$, and on
  sheets $k=q, \dots, n-1$ ($\beta$-sector) we insert $V_{\beta}$ at $t_{k}^{+}$
  and $V_{-\beta}$ at $t_{k}^{-}$.

  \subsection*{Special case}

  For simplicity, focus on the special case $y=1/2$ so that we need to compute
  \begin{eqnarray}
    2^{-(n-1)}\frac{\tr(\rho_{0}(\rho_{0}+\rho_{\beta})^{n-1})}{\tr(\rho_{0}^{n})}&&=2^{-(n-1)}\lb
    1+(n-1)\frac{\tr(\rho_{0}^{n-1}\rho_{\beta})}{\tr(\rho_{0}^{n})}+(n-1)\sum_{k=1}^{n-2}\frac{\tr(\rho_{0}^{k}\rho_{\beta}\rho_{0}^{n-1-k}\rho_{\beta})}{\tr(\rho_{0}^{n})}
    \right.\nn\\
    &&\left.+(n-1)\sum_{k_1=1}^{n-3}\sum_{k_2=0}^{n-3-k_1}\frac{\tr(\rho_{0}^{k_1}\rho_{\beta}\rho_{0}^{k_2}\rho_{\beta}\rho_{0}^{n-3-k_1-k_2}\rho_{\beta})}{\tr(\rho_{0}^{n})}+\cdots\rb
  \end{eqnarray}
  In this case, the first term $p=1$ gives
  \begin{eqnarray}
    &&\frac{\tr(\rho_{0}^{n-1}\rho_{\beta})}{\tr(\rho_{0}^{n})}=f(x)^{-\beta^2}\nn\\
    &&f(x)\equiv (2\sin(\pi x/n)),
  \end{eqnarray}
  whereas the $p=2$ becomes
  \begin{eqnarray}
    &&\sum_{k=1}^{n-2}\frac{\tr(\rho_{0}^{k}\rho_{\beta}\rho_{0}^{n-1-k}\rho_{\beta})}{\tr(\rho_{0}^{n})}=f_1(x)^{-2\beta^2}
    g_2(x;\beta)\nn\\
    &&g_2(x;\beta) = \sum_{k=1}^{n-2} \left[
    \frac{\sin^{2}\left(\frac{\pi k}{n}\right)}{\sin\left(\frac{\pi(k+x)}{n}\right)
    \sin\left(\frac{\pi(k-x)}{n}\right)}
    \right]^{\beta^2}\ .
  \end{eqnarray}
  For the $p=3$ term, we find
  \begin{eqnarray}
    \sum_{k_1=1}^{n-3}\sum_{k_2=0}^{n-3-k_1}\frac{\tr(\rho_{0}^{k_1}\rho_{\beta}\rho_{0}^{k_2}\rho_{\beta}\rho_{0}^{n-3-k_1-k_2}\rho_{\beta})}{\tr(\rho_{0}^{n})}=\cdots
  \end{eqnarray}

  \subsection*{Calculation of the correlator}

  The correlation function for primary vertex operators $V_{Q}= :e^{i Q \phi}:$
  is given by the product of distances raised to the power of the product of
  their charges. The chordal distance on the strip is $d(t_{i}, t_{j}) = 2 \sin(\frac{t_{i}-
  t_{j}}{2})$.
  \begin{equation}
    \langle \prod_{i}V_{Q_i}(t_{i}) \rangle \sim \prod_{i<j}\left[ 2 \sin\left( \frac{t_{i}-
    t_{j}}{2}\right) \right]^{Q_i Q_j}
  \end{equation}
  \begin{itemize}
    \item \textbf{Like charges} ($Q_{i}Q_{j}> 0$): The operators repel (regular
      zero), leading to a positive exponent $+Q_{i}Q_{j}$.

    \item \textbf{Opposite charges} ($Q_{i}Q_{j}< 0$): The operators attract (singular
      pole), leading to a negative exponent $-|Q_{i}Q_{j}|$.
  \end{itemize}

  We have the following contributions:
  \paragraph{1. Nearest Neighbor Interactions (Same sheet, same sector)}
  These are the contractions between $t_{k}^{+}$ and $t_{k}^{-}$ on the same sheet.
  \begin{itemize}
    \item $\alpha$-sector ($q$ pairs): $(\alpha, -\alpha)$ separated by
      $\Delta t = 2\pi x/n$.
      \[
        \implies \left( 2 \sin \frac{\pi x}{n}\right)^{-q \alpha^2}
      \]

    \item $\beta$-sector ($p$ pairs): $(\beta, -\beta)$ separated by
      $\Delta t = 2\pi x/n$.
      \[
        \implies \left( 2 \sin \frac{\pi x}{n}\right)^{-p \beta^2}
      \]
  \end{itemize}
  Putting the two together, we have
  \begin{eqnarray}
    \label{Nfunc} \mathcal{N}(q,p;x)=\left( 2 \sin \frac{\pi x}{n} \right)^{-(q\alpha^2+p\beta^2)}
  \end{eqnarray}

  \paragraph{2. Intra-sector Interactions (Different sheets, same sector)}
  For the $\alpha$-sector, we sum over sheet separations $k = 1, \dots, q-1$. The
  number of pairs at distance $k$ is $(q-k)$.
  \begin{itemize}
    \item Like charges $(\alpha, \alpha)$ and $(-\alpha, -\alpha)$:
      \[
        \text{Dist}= \frac{2\pi k}{n}\implies \text{Exp}= +\alpha^{2}\times 2
      \]

    \item Opposite charges $(\alpha, -\alpha)$: Two diagonals
      $\frac{2\pi k}{n}\pm \frac{2\pi x}{n}$.
      \[
        \text{Dist}= \frac{2\pi (k\pm x)}{n}\implies \text{Exp}= -\alpha^{2}\text{
        (each)}
      \]
  \end{itemize}
  Collecting these terms for the $\alpha$ sector gives $[g(q;x)]^{\alpha^2}$:
  \begin{eqnarray}
    g(q;x) &= \prod_{k=1}^{q-1} \left(
    \frac{\sin^{2}(\frac{\pi k}{n})}{\sin(\frac{\pi(k+x)}{n}) \sin(\frac{\pi(k-x)}{n})}
    \right)^{q-k} \ .
  \end{eqnarray}
  The $\beta$ sector is identical with $q \to p$ and $\alpha \to \beta$, and we
  have the overall contribution
  \begin{eqnarray}
    g(q;x)^{\alpha^2}g(p;x)^{\beta^2}\ .
  \end{eqnarray}

  \paragraph{3. Cross-sector Interactions}
  Interactions between the $q$ $\alpha$-sheets and $p$ $\beta$-sheets yield
  terms involving the product $\alpha \beta$. These collect into $[h(p,q;x)]^{\alpha\beta}$:
  \begin{align}
    \label{hfunct}h(q,p;x) & = \prod_{l=1}^{p}\prod_{k=1}^{q}\frac{\sin^{2}(\frac{\pi(k+l-1)}{n})}{\sin(\frac{\pi(k+l-1+x)}{n}) \sin(\frac{\pi(k+l-1-x)}{n})}
  \end{align}

  Putting all the pieces together, we find that the full correlator in (\ref{fullcorr})
  is
  \begin{equation}
    \label{eq:mixed_correlator}\langle \mathcal{O}_{\alpha}^{(q)}\mathcal{O}_{\beta}
    ^{(p)}\rangle = \mathcal{N}(n-p,p;x) \left[ g(n-p;x) \right]^{\alpha^2}\left[
    g(p;x) \right]^{\beta^2}\left[ h(n-p,p;x) \right]^{\alpha\beta}\ .
  \end{equation}

  We need to continue in $n$ analytically:
  \begin{eqnarray}
    \label{Fn} F(n)\equiv \text{tr}(\rho(\rho+\sigma)^{n-1}) &&= \sum_{p=0}^{n-1}
    \binom{n-1}{p} \langle \mathcal{O}_{\alpha}^{(q)} \mathcal{O}_{\beta}^{(p)} \rangle
  \end{eqnarray}

  \subsection{Analytic continuation in $n$}

  We consider the relative entropy $S(\rho_{\alpha}|| \sigma)$ where
  $\sigma = y \rho_{\alpha}+ (1-y)\rho_{\beta}$ for excited coherent states in a
  1+1D free boson CFT. The replica trick requires evaluating:
  \begin{equation}
    \text{tr}(\rho_{\alpha}\sigma^{n-1}) = \sum_{p=0}^{n-1}\binom{n-1}{p}y^{n-1-p}
    (1-y)^{p}C_{n-p, p}
  \end{equation}
  where
  $C_{n-p, p}= \langle \mathcal{O}_{\alpha}^{(n-p)}\mathcal{O}_{\beta}^{(p)}\rangle$
  is the mixed correlator on the $n$-sheeted replica surface.

  \paragraph{Step 1 Log-Sum expansion of the correlator:}
  ,{k,1,The correlator is a product of geometric functions $g$ and $h$. We focus
  on the intra-sector term $\log g(p;x)$. For integer $p$, it is defined as:
  \begin{equation}
    \log g(p;x) = \sum_{k=1}^{p-1}(p-k) \left[ 2 \log d\left(\frac{k}{n}\right) -
    \log d\left(\frac{k+x}{n}\right) - \log d\left(\frac{k-x}{n}\right) \right]
  \end{equation}
  where $d(y) = 2\sin(\pi y)$ is the distance between operator insertions. We
  utilize the Fourier expansion for $\log \sin(\pi y)$:
  \begin{equation}
    \log \sin(\pi y) = -\log 2 - \sum_{m=1}^{\infty}\frac{\cos(2\pi m y)}{m}
  \end{equation}
  Substituting this into the potential $\mathcal{V}(k,x)$ inside the sum, the
  constants cancel, yielding:
  \begin{equation}
    \mathcal{V}(k,x) = -2 \sum_{m=1}^{\infty}\frac{1}{m}\left[ 1 - \cos\left(\frac{2\pi
    m x}{n}\right) \right] \cos\left(\frac{2\pi m k}{n}\right)
  \end{equation}
  The expression above has the advantage that it has separated the contribution
  of $k$ from those of $x$, and we can now perform the sum over $k$.

  \paragraph{Step 2 Sum over $k$:}
  Following Section 6 of Calabrese, Cardy, and Tonni (2011), we perform the sum over
  $k$ for each mode $m$:
  \begin{equation}
    \log g(p;x) = -2 \sum_{m=1}^{\infty}\frac{1 - \cos(2\pi m x/n)}{m}\sum_{k=1}^{p-1}
    (p-k) \cos\left(\frac{2\pi m k}{n}\right)
  \end{equation}
  The inner sum $S_{m}(p) = \sum_{k=1}^{p-1}(p-k) \cos(k\theta_{m})$ is
  evaluated using geometric series. For integer $p$ and $\theta_{m}=2\pi m/n$:
  \begin{equation}
    S_{m}(p) = \text{Re}\left[ \frac{p e^{i\theta_m}}{1-e^{i\theta_m}}- \frac{e^{i\theta_m}(1-e^{ip\theta_m})}{(1-e^{i\theta_m})^{2}}
    \right]=\frac{-p}{2}+\frac{(1-\cos(p\theta_{m}))}{4\sin^{2}(\theta_{m}/2)},
  \end{equation}
  Therefore, we find
  \begin{eqnarray}
    \log g(p;x)=\sum_{m=1}^\infty \frac{1-\cos(\frac{2\pi m x}{n})}{m}\lb p-\frac{\sin^{2}(\pi
    p m/n)}{\sin^{2}(\pi m/n)}\rb
  \end{eqnarray}

  Define the mode kernels
  \begin{equation}
    L(n,x):=\sum_{m=1}^{\infty}\frac{1-\cos(2\pi m x/n)}{m}, \qquad K_{m}(n,x):=
    \frac{1-\cos(2\pi m x/n)}{2m\,\sin^{2}(\pi m/n)}. \label{eq:L-Km}
  \end{equation}
  to write
  \begin{equation}
    \log g(p;x)= p\,L(n,x)\;-\;\sum_{m=1}^{\infty}K_{m}(n,x)\,\bigl(1-\cos(p\theta
    _{m})\bigr). \label{eq:logg-final}
  \end{equation}
  All $p$-independent pieces can be absorbed into $\Omega_{n}^{\rm (rest)}(x)$
  later.

  A technical note to keep in mind is that to make every rearrangement fully
  justified, one may define all mode sums with a common Abel factor $e^{-\eps m}$
  in the summand and take $\eps\to 0^{+}$ at the end. This guarantees absolute
  and uniform convergence in intermediate steps.

  Next, we move to the function $h(q,p;x)$ in (\ref{hfunct}) which includes a double
  product over $(k,l)$ with $1\le k\le q$, $1\le l\le p$. Taking logs and using the
  same log-sum representation as in we did for the $g$ function yields
  \begin{equation}
    \log h(q,p;x)=\sum_{r=1}^{n-1}w_{r}(p,q)\,V(r,x), \label{eq:logh-weighted}
  \end{equation}
  where $w_{r}(p,q)$ counts the number of pairs $(k,l)$ such that $k+l-1=r$.

  It is convenient to evaluate the weighted cosine sum via generating functions.
  Let $\theta$ be arbitrary and note
  \begin{equation}
    \sum_{r=1}^{n-1}w_{r}(p,q)e^{\ii r\theta}=\sum_{k=1}^{q}\sum_{l=1}^{p}e^{\ii(k+l-1)\theta}
    =e^{-\ii\theta}\Big(\sum_{k=1}^{q}e^{\ii k\theta}\Big)\Big(\sum_{l=1}^{p}e^{\ii
    l\theta}\Big). \label{eq:w-gen}
  \end{equation}
  Now specialize to the case $q=n-p$ and $\theta=\theta_{m}=2\pi m/n$. Using $\sum
  _{k=1}^{n-p}e^{\ii k\theta_m}=-e^{-\ii p\theta_m/2}\,\frac{\sin(p\theta_{m}/2)}{\sin(\theta_{m}/2)}$
  and $\sum_{l=1}^{p}e^{\ii l\theta_m}=e^{\ii(p+1)\theta_m/2}\,\frac{\sin(p\theta_{m}/2)}{\sin(\theta_{m}/2)}$,
  the right-hand side of \eqref{eq:w-gen} becomes real and equals
  \begin{equation}
    \sum_{r=1}^{n-1}w_{r}(p,n-p)\cos(r\theta_{m}) = -\left(\frac{\sin(p\theta_{m}/2)}{\sin(\theta_{m}/2)}
    \right)^{2}= -\frac{1-\cos(p\theta_{m})}{2\sin^{2}(\theta_{m}/2)}. \label{eq:wcos-eval}
  \end{equation}
  Substituting \eqref{eq:wcos-eval} into \eqref{eq:logh-weighted} gives the exact
  mode form
  \begin{equation}
    \log h(n-p,p;x)= 2\sum_{m=1}^{\infty}K_{m}(n,x)\,\bigl(1-\cos(p\theta_{m})\bigr
    ), \label{eq:logh-final}
  \end{equation}
  with the \emph{same} $K_{m}(n,x)$ as in \eqref{eq:L-Km}.

  We are now ready to put together all the pieces back into $C_{n-p,p}$. Using
  \eqref{eq:logg-final} and \eqref{eq:logh-final}, we find
  \begin{align}
    \alpha^{2}\log g(q;x)+\beta^{2}\log g(p;x)+\alpha\beta\log h(q,p;x) & =(\alpha^{2}q+\beta^{2}p)\,L(n,x) \nonumber                                                             \\
                                                                        & \quad-(\alpha-\beta)^{2}\sum_{m=1}^{\infty}K_{m}(n,x)\,\bigl(1-\cos(p\theta_{m})\bigr) \label{eq:combo}
  \end{align}
  where we have used a similar log-sum expansion for the function $\mathcal{N}$.
  Absorbing all $p$-independent constants (including the $-(\alpha-\beta)^{2}\sum
  _{m}K_{m}$ term) into $\Omega_{n}(x)$, we obtain
  \begin{equation}
    C_{n-p,p}(x) =\Omega_{n}(x)\,\Lambda_{\alpha}^{\,n-p}\,\Lambda_{\beta}^{\,p}\,
    \exp\!\left[\sum_{m=1}^{\infty}a_{m}(n,x)\cos(p\theta_{m})\right], \label{eq:C-expcos}
  \end{equation}
  where
  \begin{equation}
    a_{m}(n,x)=(\alpha-\beta)^{2}\,K_{m}(n,x) =(\alpha-\beta)^{2}\,\frac{1-\cos(2\pi
    m x/n)}{2m\,\sin^{2}(\pi m/n)}\ . \label{eq:am}
  \end{equation}

  \subsection{Fourier/Bessel expansion and definition of the $R$-modes}
  Let $\theta:=\frac{2\pi p}{n}$ so that $\cos(p\theta_{m})=\cos(m\theta)$. Using
  $e^{a\cos(m\theta)}=\sum_{r\in\mathbb{Z}}I_{r}(a)\,e^{\ii r m\theta}$ (modified
  Bessel functions), we expand
  \begin{equation}
    \exp\!\left[\sum_{m\ge 1}a_{m}\cos(m\theta)\right] =\sum_{R\in\mathbb{Z}}W_{R}
    (n,x;\alpha,\beta)\,e^{\ii R\theta}, \label{eq:WR-Fourier}
  \end{equation}
  where the coefficients are the (infinite) Bessel convolution
  \begin{equation}
    W_{R}(n,x;\alpha,\beta)= \sum_{\{r_m\}\,:\,\sum_{m\ge1} m r_m=R}\;\prod_{m\ge1}
    I_{r_m}\!\big(a_{m}(n,x)\big), \qquad W_{-R}=W_{R}. \label{eq:WR-def}
  \end{equation}
  This allows us to rewrite
  \begin{equation}
    C_{n-p,p}(x) =\Omega_{n}\,\Lambda_{\alpha}^{\,n-p}\Lambda_{\beta}^{\,p}\; \sum
    _{R\in\mathbb{Z}}W_{R}(n)\,e^{\ii \frac{2\pi R}{n}p}. \label{eq:C-R}
  \end{equation}

  Note that with an Abel regulator in the defining mode sums, the function of
  $\theta$ in \eqref{eq:WR-Fourier} is real-analytic in a complex strip, which
  implies exponential decay of $W_{R}$ and in particular
  $\sum_{R}|W_{R}|<\infty$.

  The expression in (\ref{eq:C-R}) has the advantage that if we can swap the sum
  over $R$ with the sum over $p$, we can perform the $p$-sum explicitly. That is
  our next step. {\bf I believe swapping the two sums can be justified using Abel regulators but I am not entirely sure!}

  Plugging these back in \ref{Fn} and swapping the order of summation, we obtain
  \begin{align}
    F(n) & =\Omega_{n}\,\Lambda_{\alpha}\sum_{R\in\mathbb{Z}}W_{R}(n) \sum_{p=0}^{n-1}\binom{n-1}{p}\, (y\Lambda_{\alpha})^{n-1-p}\, \bigl((1-y)\Lambda_{\beta}e^{\ii2\pi R/n}\bigr)^{p}\nonumber \\
         & = \Omega_{n}\,\Lambda_{\alpha}\sum_{R\in\mathbb{Z}}W_{R}(n)\, \Bigl(y\Lambda_{\alpha}+(1-y)\Lambda_{\beta}e^{\ii2\pi R/n}\Bigr)^{n-1}\ . \label{eq:F-R}
  \end{align}

  Next, we would like to argue that only the $R=0$ term in the expression above contrinutes
  to the $n\to 1$ limit. To establish this, define $\eps=n-1$ and
  \begin{equation}
    A_{R}(n):=y\Lambda_{\alpha}+(1-y)\Lambda_{\beta}e^{\ii2\pi R/n},\qquad A_{0}:
    =y\Lambda_{\alpha}+(1-y)\Lambda_{\beta}>0.
  \end{equation}
  Then \eqref{eq:F-R} reads
  \begin{equation}
    F(n)=\Omega_{n}\Lambda_{\alpha}\sum_{R\in\mathbb{Z}}W_{R}(n)\,A_{R}(n)^{\eps}
    . \label{eq:F-eps}
  \end{equation}

  With an Abel regulator $W_{R}$ decays exponentially and has finite moments.
  Expand the phase for $n=1+\eps$:
  \begin{equation}
    e^{\ii2\pi R/n}=e^{\ii2\pi R/(1+\eps)}=\exp\!\bigl(-\ii2\pi R\eps+O(R\eps^{2}
    )\bigr),
  \end{equation}
  so $A_{R}(n)=A_{0}+\delta_{R}$ with $\delta_{R}=O(R\eps)$ for fixed $R$ and $\delta
  _{-R}=\overline{\delta_R}$. Using
  \begin{equation}
    A_{R}^{\eps}=\exp\!\bigl(\eps\log(A_{0}+\delta_{R})\bigr)=A_{0}^{\eps}\left[1
    +\eps\frac{\delta_{R}}{A_{0}}+O(\eps|\delta_{R}|^{2})\right],
  \end{equation}
  and the evenness $W_{-R}=W_{R}$, the potentially dangerous linear term $\sum_{R}
  W_{R}\,\delta_{R}$ cancels at $O(\eps)$ (its leading part is odd/purely
  imaginary), leaving only $O(\eps^{2})$ corrections after summing over $R$. Consequently,
  after absorbing $\sum_{R}W_{R}$ into $\Omega_{n}$ (a $p$-independent normalization),
  \begin{equation}
    F(1+\eps)=\Omega_{1+\eps}\Lambda_{\alpha}\,A_{0}^{\eps}\bigl(1+O(\eps^{2})\bigr
    ). \label{eq:F-asymp}
  \end{equation}
  Thus the $R\neq 0$ modes do not affect $\partial_{n}\log F(n)$ at $n=1$.

  The replica prescription for the relative entropy is
  \begin{equation}
    S(\rho_{\alpha}\|\sigma_{y})= \left.\partial_{n}\Bigl(\log \tr(\rho_{\alpha}^{n}
    )-\log \tr(\rho_{\alpha}\sigma_{y}^{n-1})\Bigr)\right|_{n=1}. \label{eq:S-replica}
  \end{equation}
  In the same framework one has $\tr(\rho_{\alpha}^{n})=\widetilde\Omega_{n}\,\Lambda
  _{\alpha}^{\,n}$ so that $\partial_{n}\log \tr(\rho_{\alpha}^{n})|_{n=1}=\log\Lambda
  _{\alpha}$ (all $\widetilde \Omega_{n}$ terms cancel against $\Omega_{n}$ in
  the difference). Using \eqref{eq:F-asymp},
  \begin{equation}
    \left.\partial_{n}\log \tr(\rho_{\alpha}\sigma_{y}^{n-1})\right|_{n=1}=\log A
    _{0}=\log\!\bigl(y\Lambda_{\alpha}+(1-y)\Lambda_{\beta}\bigr).
  \end{equation}
  Therefore
  \begin{equation}
    S(\rho_{\alpha}\|\sigma_{y})=\log\Lambda_{\alpha}-\log\!\bigl(y\Lambda_{\alpha}
    +(1-y)\Lambda_{\beta}\bigr) = -\log\!\left[y+(1-y)\frac{\Lambda_{\beta}}{\Lambda_{\alpha}}
    \right]. \label{eq:S-Lambda}
  \end{equation}

  For coherent-state excitations in the free boson CFT one may identify
  \begin{equation}
    \frac{\Lambda_{\beta}}{\Lambda_{\alpha}}=e^{-S(\alpha\|\beta)}, \qquad S(\alpha
    \|\beta)=(\alpha-\beta)^{2}K(x), \qquad K(x)=1-\pi x\cot(\pi x). \label{eq:Salpha}
  \end{equation}
  Substituting \eqref{eq:Salpha} into \eqref{eq:S-Lambda} gives the final answer
  \begin{equation}
    \boxed{ S\!\left(\rho_\alpha\,\Big\|\,y\rho_\alpha+(1-y)\rho_\beta\right) = -\log\!\left( y+(1-y)\exp\!\big[-(\alpha-\beta)^2(1-\pi x\cot(\pi x))\big] \right). }
    \label{eq:S-final}
  \end{equation}

  \section{Consistency tests}
  \label{sec:tests}

  Let $s:=(\alpha-\beta)^{2}K(x)\ge 0$ and $a:=e^{-s}\in(0,1]$ so that
  \begin{eqnarray}
    S(y)\equiv S\!\left(\rho_\alpha\,\Big\|\,y\rho_\alpha+(1-y)\rho_\beta\right)=-\log\lb
    y+(1-y)a\rb=-\log(a+(1-y)a)
  \end{eqnarray}
  We are gonna test the following limits:
  \begin{itemize}
    \item \textbf{Identical states:}
      $\alpha=\beta\Rightarrow s=0\Rightarrow a=1\Rightarrow S(y)=0$ as expected.

    \item \textbf{No mixing:} $S(1)\to 0$ as expected and
      $S(0)=S(\alpha\|\beta)$ as expected. Therefore, we can write
      \begin{eqnarray}
        e^{-S(y)}=y e^{-S(1)}+(1-y)e^{-S(0)}\ .
      \end{eqnarray}

    \item \textbf{Much heavier state:}
      $s\to\infty\Rightarrow a\to 0\Rightarrow S_{y}\to -\log y$.
  \end{itemize}

  Next, we test the following properties:

  \begin{itemize}
    \item \textbf{Positivity:} Since $y\le a+y(1-a)\le 1$, we have $0\le S(y)\le
      -\log y$ as required fo r relative entropy.

    \item \textbf{Convexity in $y$:} Taking derivatives with respect to $y$
      \begin{equation}
        S'(y)= -\frac{1-a}{a+y(1-a)}\le 0, \qquad S''(y)=\frac{(1-a)^{2}}{(a+y(1-a))^{2}}
        \ge 0.
      \end{equation}
      Thus $S(y)$ is decreasing and convex in the mixing parameter $y$ (strictly
      convex if $s>0$), consistent with the general theorem that $S(\rho\|\sigma)$
      is convex in $\sigma$.

    \item \textbf{Monotonicity and concavity in $s$:} Treating $S$ as a function
      of $s$:
      \begin{equation}
        \frac{\partial S}{\partial s}=\frac{(1-y)e^{-s}}{y+(1-y)e^{-s}}\ge 0, \qquad
        \frac{\partial^{2}S}{\partial s^{2}}=-\frac{y(1-y)e^{-s}}{(y+(1-y)e^{-s})^{2}}
        \le 0.
      \end{equation}
      So $S$ increases with $s$ but is concave in $s$, consistent with
      saturation at $-\log y$. A direct derivative gives the derivative of relative
      entropy in region size
      \begin{equation}
        K'(x)=\frac{\pi}{\sin^{2}(\pi x)}\left(\pi x-\frac{1}{2}\sin(2\pi x)\right
        )>0 \qquad \text{for }x\in(0,1),
      \end{equation}
      since $\sin u<u$ for $u>0$. Thus $K(x)$ (and therefore $S$) is increasing
      in $x$ for $\alpha\neq\beta$, as expected from data processing.
  \end{itemize}

  Finally, we note that in the small interval limit $x\ll 1$ we have
  \begin{equation}
    K(x)=1-\pi x\cot(\pi x)=\frac{\pi^{2}}{3}x^{2}+\frac{\pi^{4}}{45}x^{4}+O(x^{6}
    ).
  \end{equation}
  Hence
  \begin{equation}
    S(y)= (1-y)\,(\alpha-\beta)^{2}\frac{\pi^{2}}{3}x^{2}+O(x^{4}).
  \end{equation}
  The relative entropy vanishes as $x^{2}$ for small intervals, as expected.

  \subsection*{Alternative analytic continuation}
  \label{sec:small-n1} An alternative analytic continuation follows from observing
  that since we are interested in small $n-1$, and the sum over
  $p=1,\cdots , n-1$, we linearize our functions in $p$ and discard higher order
  terms:
  \begin{equation}
    S_{m}(p)= -\frac{p}{2}+\frac{1-\cos(p\theta_{m})}{4\sin^{2}(\theta_{m}/2)}= -
    \frac{p}{2}+O(p^{2}). \label{eq:Sm-smallp}
  \end{equation}
  Expanding the $g$ function in $n=1+\epsilon$ results in
  \begin{equation}
    \log g(p;x)=p\,L(n,x)+O(p^{2})=p\,L(n,x)+O(\eps^{2}),
  \end{equation}
  uniformly over the relevant range $p\in\{0,\dots,n-1\}$. This reproduces the linear-in-$p$
  structure that one would obtain by keeping only the ``smooth'' term $S_{m}(p)\approx
  -p/2$ and effectively discarding the oscillatory $\cos(p\theta_{m})$ piece at this
  order. While this results in a different analytic continuation than what we
  computed above, the $n\to 1$ limit gives the same result, and the same final
  expression \eqref{eq:S-final} for relative entropy.

  %------------------------------------------------------------
  % Appendix-style writeup: pole-matching analytic continuation
  %------------------------------------------------------------

  \section*{Pole-matching analytic continuation of $g(p;x)$ and $h(n-p,p;x)$}

  Here, we reformulate the analytic continuation of the building blocks $g(p;x)$
  and $h(n-p,p;x)$ using the method used in the appendix of \texttt{arXiv:1508.03506}.
  The idea is to this: To find the analytic cotinuation of a meromorphic function
  $f_{1}(z)$, you find another entire function with the same poles and zeros
  $f_{2}(z)$ such that $f_{1}(z)/f_{z}(2)$ is bounded. Then, by Liouville's
  theorem it must be a constant $C$. This means that $C f_{2}(z)$ is the analytic
  continuation of $f_{1}(z)$.

  \subsection*{The basic building block and the definitions}

  Introduce
  \begin{equation}
    f_{n,m}(x) :=\frac{\sin^{2}(\pi m/n)}{\sin\!\bigl(\pi(m+x)/n\bigr)\,\sin\!\bigl(\pi(m-x)/n\bigr)}
    \,, \qquad m\in\mathbb{Z}.
  \end{equation}
  For integer $n\ge 2$ and $p\in\{1,\dots,n-1\}$, the functions appearing in the
  correlator are
  \begin{align}
    g(p;x)   & =\prod_{m=1}^{p-1}f_{n,m}(x)^{\,p-m}, \label{eq:g-def-poleapp}                                                                  \\[2pt]
    h(q,p;x) & =\prod_{l=1}^{p}\prod_{k=1}^{q}f_{n,k+l-1}(x) =\prod_{r=1}^{n-1}f_{n,r}(x)^{\,w_r(p,q)}, \qquad q=n-p, \label{eq:h-def-poleapp}
  \end{align}
  where $w_{r}(p,q)$ counts the number of pairs $(k,l)$ such that $k+l-1=r$.

  \subsection*{Poles of $g(p;x)$ as a function of $x$}

  From \eqref{eq:g-def-poleapp}, all $x$-dependence sits in the denominator
  sines. For fixed $m$, the factor $\sin(\pi(m\pm x)/n)$ vanishes iff
  $m\pm x=\ell n$ for some $\ell\in\mathbb{Z}$, i.e.
  \begin{equation}
    x=\pm m+\ell n.
  \end{equation}
  Hence the pole set of $g(p;x)$ is
  \begin{equation}
    \boxed{ \mathrm{Pole}(g)=\Bigl\{\,x=\pm m+\ell n\ \Big|\ m=1,\dots,p-1,\ \ell\in\mathbb{Z}\Bigr\}. }
  \end{equation}
  Moreover, each vanishing sine gives a simple pole, and the factor $f_{n,m}(x)$
  is raised to power $(p-m)$. Therefore the \emph{pole order} of $g(p;x)$ at
  $x=\pm m+\ell n$ equals $(p-m)$:
  \begin{equation}
    \boxed{ \mathrm{ord}_{x=\pm m+\ell n}\, g(p;x)=p-m. }
  \end{equation}
  Finally, $g(p;x)$ has no zeros as a function of $x$, since the numerator $\sin^{2}
  (\pi m/n)$ is $x$-independent.

  \subsection*{A multiple-sine candidate and its divisor}

  We now introduce a candidate $g_{\rm MS}(p;x)$ built from multiple sine
  functions. Let $S_{2}(z|1,n)$ and $S_{3}(z|1,1,n)$ denote the double and triple
  sine functions (in a standard normalization), characterized by the shift
  relations
  \begin{equation}
    \frac{S_{2}(z+1\,|\,1,n)}{S_{2}(z\,|\,1,n)}=2\sin\!\Bigl(\frac{\pi z}{n}\Bigr
    ), \qquad \frac{S_{3}(z+1\,|\,1,1,n)}{S_{3}(z\,|\,1,1,n)}=\frac{1}{S_{2}(z\,|\,1,n)}
    . \label{eq:S2S3-shifts}
  \end{equation}
  Define the triangular-product package
  \begin{equation}
    T_{p}(z):= \frac{S_{3}(z+2\,|\,1,1,n)^{\,p}}{S_{3}(z+p+1\,|\,1,1,n)\,S_{3}(z+1\,|\,1,1,n)^{\,p-1}}
    , \label{eq:T-def}
  \end{equation}
  and the candidate analytic continuation
  \begin{equation}
    g_{\rm MS}(p;x):=\frac{T_{p}(0)^{2}}{T_{p}(x)\,T_{p}(-x)}. \label{eq:gMS-def}
  \end{equation}

  \paragraph{Claim (telescoping to the finite sine product).}
  By iterating the shift relations \eqref{eq:S2S3-shifts}, one checks that
  $T_{p}(z)$ is proportional to the triangular sine product
  \begin{equation}
    T_{p}(z)\ \propto\ \prod_{m=1}^{p-1}\sin\!\Bigl(\frac{\pi(m+z)}{n}\Bigr)^{p-m}
    , \label{eq:T-telescope}
  \end{equation}
  where the proportionality factor is independent of $z$ (it may depend on $n,p$).
  Consequently, \eqref{eq:gMS-def} implies that for integer $n$,
  $g_{\rm MS}(p;x)$ reproduces the same pole lattice and the same pole orders as
  $g(p;x)$.

  \paragraph{Zeros/poles of $g_{\rm MS}(p;x)$.}
  From \eqref{eq:T-telescope}, $T_{p}(z)$ has zeros whenever $z=\ell n-m$ for $m=
  1,\dots,p-1$, with zero order $(p-m)$. Therefore $g_{\rm MS}(p;x)$ has poles whenever
  $T_{p}(x)=0$ or $T_{p}(-x)=0$, i.e.
  \begin{equation}
    x=\ell n-m\quad \text{or}\quad x=-(\ell n-m)=m-\ell n \quad\Longleftrightarrow
    \quad x=\pm m+\ell n,
  \end{equation}
  and the pole order at each such point equals $(p-m)$. In particular,
  \begin{equation}
    \boxed{ \mathrm{Pole}(g_{\rm MS})=\mathrm{Pole}(g),\qquad \mathrm{ord}_{x=\pm m+\ell n}\,g_{\rm MS}(p;x)=p-m. }
  \end{equation}
  Like $g$, the function $g_{\rm MS}$ has no zeros as a function of $x$.

  \subsection*{The ratio has no poles and is bounded}

  Define the ratio
  \begin{equation}
    R(x):=\frac{g_{\rm MS}(p;x)}{g(p;x)}.
  \end{equation}
  Since $g_{\rm MS}$ and $g$ have the same poles with the same orders and no zeros,
  all singularities cancel:
  \begin{equation}
    \boxed{R(x)\ \text{is entire in }x.}
  \end{equation}
  Moreover, both $g$ and $g_{\rm MS}$ are even and $n$-periodic in $x$, hence so
  is $R$:
  \begin{equation}
    R(-x)=R(x),\qquad R(x+n)=R(x). \label{eq:R-periodic}
  \end{equation}

  To argue boundedness, it suffices (by periodicity) to bound $R$ on a single
  vertical strip $\{x=u+\ii v: u\in[0,n]\}$. On such a strip, the original
  $g(p;x)$ obeys an exponential-type estimate as $|v|\to\infty$: for each fixed $m$,
  \begin{equation}
    \Bigl|\sin\!\Bigl(\frac{\pi(m\pm(u+\ii v))}{n}\Bigr)\Bigr| \sim \frac{1}{2}e^{\pi
    |v|/n}, \qquad |v|\to\infty,
  \end{equation}
  so
  \begin{equation}
    |g(p;u+\ii v)| \sim C\,\exp\!\Bigl[-\frac{2\pi|v|}{n}\sum_{m=1}^{p-1}(p-m)\Bigr
    ] = C\,\exp\!\Bigl[-\frac{\pi p(p-1)}{n}|v|\Bigr], \qquad |v|\to\infty, \label{eq:g-asymp}
  \end{equation}
  with $C$ independent of $v$ (uniformly for $u$ in compact sets away from poles).

  On the other hand, the multiple sine functions $S_{2},S_{3}$ have well-known large-imaginary
  asymptotics of the form
  \begin{equation}
    \log S_{r}(u+\ii v\,|\,\boldsymbol{\omega}) = \pm \frac{\pi}{r!\,\omega_{1}\cdots
    \omega_{r}}\,|v|^{r}\ +\ \text{lower powers in }|v| \ +\ O(1), \qquad |v|\to
    \infty,
  \end{equation}
  so that in ratios such as \eqref{eq:T-def} the leading polynomial growth
  cancels, and $T_{p}(z)$ has exponential type matching the finite sine product in
  \eqref{eq:T-telescope}. Consequently $g_{\rm MS}(p;u+\ii v)$ has the same exponential
  type as $g(p;u+\ii v)$, and hence their ratio $R(u+\ii v)$ remains bounded as
  $|v|\to\infty$ uniformly in $u\in[0,n]$. Together with continuity on the compact
  region $\{u\in[0,n],\,|v|\le V\}$, we conclude:
  \begin{equation}
    \boxed{R(x)\ \text{is bounded on }\mathbb{C}.}
  \end{equation}

  Since $R$ is entire and bounded, Liouville's theorem implies $R$ is constant.
  Finally, $g(p;0)=1$ by definition and $g_{\rm MS}(p;0)=1$ by \eqref{eq:gMS-def},
  so $R(0)=1$ and thus
  \begin{equation}
    \boxed{g_{\rm MS}(p;x)=g(p;x)\quad \text{for all }x\ \text{(at integer }n\text{).}}
  \end{equation}
  This is precisely the ``same poles + bounded ratio'' strategy.

  \subsection*{Corollary: $h(n-p,p;x)$ from a full-product identity}

  Using \eqref{eq:h-def-poleapp} with $q=n-p$ and the explicit multiplicities
  \begin{equation}
    w_{r}(p,n-p)=
    \begin{cases}
      r,   & 1\le r\le p,  \\
      p,   & p<r\le n-p,   \\
      n-r, & n-p<r\le n-1,
    \end{cases}
  \end{equation}
  one may rewrite $h(n-p,p;x)$ as
  \begin{equation}
    h(n-p,p;x)=\Bigl(\prod_{r=1}^{n-1}f_{n,r}(x)\Bigr)^{p}\,g(p;x)^{-2}. \label{eq:h-via-fullprod}
  \end{equation}
  The appendix identity of \texttt{arXiv:1508.03506} proves
  \begin{equation}
    \prod_{r=1}^{n-1}f_{n,r}(x)=\left(\frac{n\sin(\pi x/n)}{\sin(\pi x)}\right)^{2}
    ,
  \end{equation}
  again by the same poles/bounded ratio argument. Substituting into \eqref{eq:h-via-fullprod}
  yields
  \begin{equation}
    \boxed{ h(n-p,p;x) = \left(\frac{n\sin(\pi x/n)}{\sin(\pi x)}\right)^{2p}\,g(p;x)^{-2}, }
  \end{equation}
  so the analytic continuation of $h$ follows once that of $g$ is fixed by the
  pole-matching method.

  \section{Information-theoretic interpretation}
\end{document}